\section{Lineamientos para Implementación Educativa}

Este capítulo establece los lineamientos para la implementación del kit PhantomX Pincher en entornos educativos, proporcionando una guía estructurada para docentes y estudiantes que utilizarán la plataforma desarrollada en actividades de laboratorio. El enfoque pedagógico propuesto permite a los estudiantes construir conocimiento progresivo en robótica, visión por computador y programación mediante un proyecto integrador.

%==============================================================================
\subsection{Objetivos de Aprendizaje}
%==============================================================================

La implementación del kit en un curso de robótica permite alcanzar los siguientes objetivos educativos:

\begin{enumerate}
    \item Configurar y gestionar entornos de desarrollo para ROS2 Jazzy en sistemas Ubuntu 24.04 LTS
    \item Comprender y modificar modelos cinemáticos de robots mediante formato URDF/Xacro
    \item Utilizar MoveIt2 para planificación de trayectorias en manipuladores seriales
    \item Implementar nodos ROS2 con comunicación mediante tópicos y servicios
    \item Integrar procesamiento de imágenes con OpenCV para detección de objetos
    \item Desarrollar rutinas de manipulación automatizadas para tareas de clasificación
    \item Aplicar conceptos de cinemática directa e inversa en manipuladores de 4 grados de libertad
    \item Documentar proyectos técnicos mediante estándares de ingeniería
\end{enumerate}

Estos objetivos alinean con competencias de programas de Ingeniería Mecatrónica, Ingeniería Electrónica e Ingeniería de Sistemas, cubriendo áreas de robótica, control automático, visión artificial y sistemas embebidos.

%==============================================================================
\subsection{Estructura Pedagógica Propuesta}
%==============================================================================

El proyecto educativo se organiza en seis fases secuenciales que permiten construcción progresiva de conocimiento, minimizando la carga cognitiva mediante modularidad y abstracción incremental.

\subsubsection{Fase 1: Configuración del Entorno de Desarrollo}

\paragraph{Objetivos Específicos}
\begin{itemize}
    \item Instalar ROS2 Jazzy Desktop en Ubuntu 24.04 LTS
    \item Configurar workspace de desarrollo con estructura estándar de ROS2
    \item Instalar dependencias de MoveIt2, control y visualización
    \item Verificar funcionamiento mediante ejemplos básicos
\end{itemize}

\paragraph{Requisitos Previos}
\begin{itemize}
    \item Sistema operativo Ubuntu 24.04 LTS (instalación nativa o máquina virtual con mínimo 4 GB RAM, 50 GB disco)
    \item Conocimientos básicos de terminal Linux (navegación, permisos, edición de archivos)
    \item Conexión a internet estable para descarga de paquetes (~5 GB)
\end{itemize}

\paragraph{Actividades}

Los estudiantes deben seguir la documentación oficial de ROS2 Jazzy (\url{https://docs.ros.org/en/jazzy/}) para instalar la distribución completa Desktop. Posteriormente, deben configurar el entorno permanentemente agregando el script de inicialización al archivo de configuración de bash (\texttt{\textasciitilde/.bashrc}).

La instalación de dependencias del sistema incluye herramientas de desarrollo (terminal avanzado, control de versiones), así como los siguientes paquetes ROS2 esenciales para el proyecto:

\begin{itemize}
    \item \textbf{Control}: \texttt{ros2-control} y \texttt{ros2-controllers} para gestión de controladores de hardware
    \item \textbf{Modelado}: \texttt{xacro} para procesamiento de macros en archivos URDF
    \item \textbf{Visualización}: \texttt{joint-state-publisher-gui} para control manual de articulaciones
    \item \textbf{Transformadas}: \texttt{tf-transformations} para cálculos de marcos de referencia
    \item \textbf{Planificación}: Suite completa \texttt{moveit} para planificación de movimiento
\end{itemize}

El workspace del proyecto se clona desde el repositorio institucional y se compila utilizando la herramienta \texttt{colcon} con opción de enlaces simbólicos para facilitar desarrollo iterativo.

\paragraph{Resultados Esperados}

Al finalizar esta fase, el estudiante debe poder:
\begin{itemize}
    \item Ejecutar \texttt{ros2 topic list} y visualizar tópicos del sistema
    \item Lanzar visualización del robot en RViz mediante \texttt{ros2 launch phantomx\_pincher\_description view.launch.py}
    \item Mover articulaciones mediante interfaz gráfica \texttt{joint\_state\_publisher\_gui}
\end{itemize}

\subsubsection{Fase 2: Modelado 3D y Descripción Cinemática}

\paragraph{Objetivos Específicos}
\begin{itemize}
    \item Comprender formato URDF y sistema de macros Xacro
    \item Modificar modelo del robot agregando componentes personalizados
    \item Integrar mallas STL de contenedor y zona de recolección
    \item Validar modelo mediante visualización en RViz
\end{itemize}

\paragraph{Actividades}

Los estudiantes analizan la estructura del archivo URDF principal, que contiene la definición del robot mediante el sistema de macros Xacro. El archivo define argumentos configurables como \texttt{ros2\_control} (habilitar/deshabilitar sistema de control), \texttt{collision} (incluir geometrías de colisión), y otros parámetros estructurales. Los componentes del robot (brazo y gripper) se incluyen mediante directivas \texttt{xacro:include} que referencian archivos separados, facilitando modularidad y reutilización.

Para agregar componentes personalizados (contenedor, zona de recolección, canecas), los estudiantes deben crear nuevos elementos \texttt{link} con geometrías tipo mesh que referencian los archivos STL previamente diseñados. Estos links se conectan al modelo mediante elementos \texttt{joint} de tipo fijo, especificando el enlace padre (\texttt{base\_link}) y las transformadas de posición y orientación relativas.

\paragraph{Resultados Esperados}

Modelo del robot que incluye:
\begin{itemize}
    \item Manipulador PhantomX Pincher (4 DOF) con mallas visuales y de colisión
    \item Contenedor de almacenamiento integrado en \texttt{base\_link}
    \item Zona de recolección circular (Ø146 mm) en posición accesible
    \item Mástil de cámara con altura ajustable
    \item Cuatro canecas de depósito en posiciones fijas
\end{itemize}

\subsubsection{Fase 3: Planificación de Movimiento con MoveIt2}

\paragraph{Objetivos Específicos}
\begin{itemize}
    \item Comprender arquitectura de MoveIt2 (grupos de planificación, escena, controladores)
    \item Utilizar RViz Motion Planning Plugin para planificación interactiva
    \item Implementar rutinas de movimiento para clasificación de objetos
    \item Analizar trayectorias mediante visualización de estados intermedios
\end{itemize}

\paragraph{Conceptos Fundamentales}

\textbf{Grupos de Planificación}: El robot define dos grupos:
\begin{itemize}
    \item \texttt{arm}: Cadena cinemática de 4-DOF (shoulder pan, shoulder lift, elbow flex, wrist flex)
    \item \texttt{gripper}: 2-DOF para apertura/cierre de dedos
\end{itemize}

\textbf{Estados Nombrados}: Configuraciones predefinidas accesibles por nombre:
\begin{itemize}
    \item \texttt{up}: Posición vertical para transporte
    \item \texttt{rest}: Posición de reposo fuera de zona de trabajo
    \item \texttt{ready\_near}: Posición preparatoria sobre zona de recolección
\end{itemize}

\paragraph{Actividades}

Los estudiantes inician la simulación completa ejecutando el archivo launch principal del paquete \texttt{phantomx\_pincher\_bringup}, que arranca todos los nodos necesarios: controladores, MoveIt2, RViz y robot state publisher.

La planificación manual en RViz sigue un flujo interactivo:
\begin{enumerate}
    \item Seleccionar grupo de planificación \texttt{arm}
    \item Mover marcador interactivo a posición objetivo
    \item Presionar \textit{Plan} para generar trayectoria
    \item Visualizar estados intermedios
    \item Presionar \textit{Execute} para ejecutar en simulación
\end{enumerate}

Los estudiantes definen la lógica de clasificación siguiendo una secuencia de 8 pasos: (1) mover a posición preparatoria sobre zona de recolección, (2) descender a altura de agarre, (3) cerrar gripper, (4) ascender a altura segura, (5) consultar coordenadas de caneca según tipo de figura, (6) mover a posición sobre caneca correspondiente, (7) abrir gripper para liberar objeto, (8) regresar a posición de reposo. Esta secuencia garantiza trayectorias libres de colisión y operación consistente.

\paragraph{Resultados Esperados}

Estudiantes capaces de:
\begin{itemize}
    \item Planificar trayectorias libres de colisiones entre posiciones arbitrarias
    \item Identificar estados singulares del manipulador (configuraciones no alcanzables)
    \item Evaluar calidad de trayectorias (suavidad, longitud, tiempo de ejecución)
    \item Implementar rutinas de clasificación para las 4 figuras geométricas
\end{itemize}

\subsubsection{Fase 4: Programación de Nodos ROS2}

\paragraph{Objetivos Específicos}
\begin{itemize}
    \item Implementar nodos ROS2 en Python con arquitectura publisher/subscriber
    \item Utilizar mensajes estándar (\texttt{std\_msgs}, \texttt{geometry\_msgs})
    \item Definir mensajes personalizados para comandos de pose
    \item Integrar nodos con sistema MoveIt2 mediante topics
\end{itemize}

\paragraph{Arquitectura de Comunicación}

\begin{table}[H]
\centering
\caption{Tópicos de comunicación del sistema}
\label{tab:topics}
\small
\begin{tabular}{|l|l|p{6cm}|}
\hline
\textbf{Tópico} & \textbf{Tipo} & \textbf{Descripción} \\
\hline
\texttt{/tipo\_figura} & \texttt{std\_msgs/String} & Publicado por nodo de visión, indica tipo de figura detectada \\
\texttt{/pose\_command} & \texttt{PoseCommand} & Publicado por clasificador, envía pose objetivo a MoveIt \\
\texttt{/joint\_states} & \texttt{sensor\_msgs/JointState} & Publicado por controladores, estados actuales de articulaciones \\
\texttt{/gripper\_command} & \texttt{std\_msgs/Float64MultiArray} & Control de apertura/cierre del gripper \\
\hline
\end{tabular}
\end{table}

\paragraph{Implementación de Nodo Clasificador}

El nodo clasificador se implementa en Python utilizando la biblioteca \texttt{rclpy}. La estructura básica incluye:

\begin{itemize}
    \item \textbf{Inicialización}: Crear nodo heredando de clase \texttt{Node}, configurar suscriptor al tópico \texttt{/tipo\_figura} con callback asociado
    \item \textbf{Diccionario de coordenadas}: Mapear cada tipo de figura (cubo, rectángulo, pentágono, círculo) a coordenadas XYZ de caneca correspondiente
    \item \textbf{Callback de figura}: Función que se ejecuta al recibir mensaje en tópico, valida tipo de figura y dispara rutina de clasificación
    \item \textbf{Ejecución de rutina}: Implementar secuencia de movimientos comunicándose con MoveIt2 mediante mensajes de pose objetivo
    \item \textbf{Main loop}: Inicializar ROS2, crear instancia del nodo, mantener nodo activo con \texttt{spin}, destruir al finalizar
\end{itemize}

Los estudiantes deben completar la función \texttt{ejecutar\_rutina} integrando llamadas a MoveIt2 para planificación y ejecución de trayectorias.

\paragraph{Resultados Esperados}

Nodo funcional que:
\begin{itemize}
    \item Se suscribe a tópico de detección de figuras
    \item Mapea tipo de figura a coordenadas de caneca correspondiente
    \item Publica comandos de movimiento a MoveIt2
    \item Registra eventos en log de ROS2
    \item Maneja excepciones y estados de error
\end{itemize}

\subsubsection{Fase 5: Visión por Computador}

\paragraph{Objetivos Específicos}
\begin{itemize}
    \item Capturar imágenes desde cámara USB mediante \texttt{cv\_bridge}
    \item Implementar pipeline de procesamiento de imágenes con OpenCV
    \item Detectar y clasificar figuras geométricas mediante análisis de contornos
    \item Publicar resultados de detección en tópico ROS2
\end{itemize}

\paragraph{Figuras y Correspondencia}

\begin{table}[H]
\centering
\caption{Especificación de figuras a detectar}
\label{tab:figuras}
\small
\begin{tabular}{|l|l|l|l|}
\hline
\textbf{Figura} & \textbf{Caneca} & \textbf{Color} & \textbf{Dimensiones} \\
\hline
Cubo & Roja & Naranja TPU & 25 × 25 × 25 mm \\
Rectángulo & Amarilla & Naranja TPU & 50 × 25 × 25 mm \\
Pentágono & Azul & Naranja TPU & Ø25 × 25 mm (inscrito) \\
Círculo & Verde & Naranja TPU & Ø25 × 25 mm \\
\hline
\end{tabular}
\end{table}

\paragraph{Pipeline de Procesamiento}

\begin{enumerate}
    \item \textbf{Captura}: Obtener frame de cámara Logitech C270 (1280×720 @ 30 fps)
    \item \textbf{Conversión}: RGB a HSV para segmentación por color naranja
    \item \textbf{Umbralización}: Crear máscara binaria de píxeles en rango HSV del TPU naranja
    \item \textbf{Morfología}: Aplicar operaciones de cierre para eliminar ruido
    \item \textbf{Detección de contornos}: Encontrar contornos cerrados mediante algoritmo de Suzuki
    \item \textbf{Filtrado}: Descartar contornos con área menor a umbral (evitar falsos positivos)
    \item \textbf{Aproximación poligonal}: Aproximar contorno mediante Douglas-Peucker
    \item \textbf{Clasificación}: Contar vértices y calcular aspect ratio para identificar figura
    \item \textbf{Publicación}: Enviar tipo de figura detectada a tópico \texttt{/tipo\_figura}
\end{enumerate}

\paragraph{Implementación de Detección}

La función de clasificación de figuras utiliza el algoritmo de aproximación poligonal de Douglas-Peucker disponible en OpenCV. El proceso incluye:

\begin{enumerate}
    \item \textbf{Aproximación poligonal}: Calcular perímetro del contorno, aplicar tolerancia del 2\% para simplificación, obtener polígono aproximado
    \item \textbf{Conteo de vértices}: Determinar número de vértices del polígono aproximado
    \item \textbf{Clasificación por forma}:
    \begin{itemize}
        \item 4 vértices: Calcular aspect ratio del rectángulo envolvente para diferenciar cubo (ratio cercano a 1.0) de rectángulo (ratio alejado de 1.0)
        \item 5 vértices: Identificar como pentágono
        \item Más de 6 vértices: Clasificar como círculo por aproximación a forma circular
    \end{itemize}
    \item \textbf{Retorno}: Devolver string con tipo de figura detectada o "desconocido" si no coincide con patrones esperados
\end{enumerate}

Los umbrales de aspect ratio (0.9 a 1.1 para cubos) deben ajustarse experimentalmente según calibración de cámara y distorsión de lente.

\paragraph{Calibración de Parámetros HSV}

Determinación experimental de rangos HSV para TPU naranja bajo iluminación artificial (400-600 lux):

\begin{itemize}
    \item \textbf{Hue}: 10-25 (naranja en espacio HSV)
    \item \textbf{Saturation}: 100-255 (alta saturación para evitar confusión con fondo blanco)
    \item \textbf{Value}: 100-255 (rango amplio para compensar variaciones de iluminación)
\end{itemize}

\paragraph{Resultados Esperados}

Sistema de visión capaz de:
\begin{itemize}
    \item Detectar las 4 figuras con precisión >90\% bajo iluminación controlada
    \item Procesar frames a mínimo 10 fps para operación en tiempo real
    \item Filtrar falsos positivos mediante umbrales de área y circularidad
    \item Publicar detecciones con latencia <100 ms
\end{itemize}

\subsubsection{Fase 6: Integración y Validación del Sistema}

\paragraph{Objetivos Específicos}
\begin{itemize}
    \item Integrar todos los componentes en sistema funcional completo
    \item Validar funcionamiento mediante pruebas de clasificación
    \item Optimizar parámetros de velocidad y precisión
    \item Documentar sistema mediante diagramas y videos demostrativos
\end{itemize}

\paragraph{Arquitectura Integrada}

\begin{figure}[H]
\centering
\begin{tikzpicture}[node distance=2cm, auto]
    \tikzstyle{node} = [rectangle, draw, fill=blue!20, text width=5em, text centered, minimum height=2.5em]
    \tikzstyle{topic} = [ellipse, draw, fill=green!20, text width=4em, text centered]
    \tikzstyle{arrow} = [thick,-latex]

    \node (cam) [node] {Cámara\\C270};
    \node (vision) [node, right of=cam, xshift=2cm] {Nodo\\Visión};
    \node (topic1) [topic, right of=vision, xshift=1.5cm] {/tipo\\figura};
    \node (clasif) [node, right of=topic1, xshift=1.5cm] {Nodo\\Clasificador};
    \node (topic2) [topic, below of=clasif] {/pose\\command};
    \node (moveit) [node, below of=topic2] {MoveIt2};
    \node (robot) [node, left of=moveit, xshift=-2cm] {Robot\\Físico};

    \draw [arrow] (cam) -- (vision);
    \draw [arrow] (vision) -- (topic1);
    \draw [arrow] (topic1) -- (clasif);
    \draw [arrow] (clasif) -- (topic2);
    \draw [arrow] (topic2) -- (moveit);
    \draw [arrow] (moveit) -- (robot);
\end{tikzpicture}
\caption{Arquitectura del sistema de clasificación integrado. El flujo de información va desde la captura visual hasta la ejecución física, pasando por procesamiento, planificación y control.}
\label{fig:arquitectura_sistema}
\end{figure}

\paragraph{Protocolo de Pruebas}

\begin{enumerate}
    \item \textbf{Prueba de Detección Visual}
    \begin{itemize}
        \item Colocar cada figura individualmente en zona de recolección
        \item Verificar detección correcta en tópico \texttt{/tipo\_figura}
        \item Medir tasa de aciertos sobre 20 intentos por figura
    \end{itemize}
    
    \item \textbf{Prueba de Rutinas de Clasificación}
    \begin{itemize}
        \item Ejecutar rutina para cada tipo de figura
        \item Verificar trayectoria libre de colisiones
        \item Confirmar depósito en caneca correcta
    \end{itemize}
    
    \item \textbf{Prueba de Sistema Completo}
    \begin{itemize}
        \item Secuencia aleatoria de 16 figuras (4 de cada tipo)
        \item Medir tiempo total de clasificación
        \item Registrar errores de detección o manipulación
        \item Calcular tasa de éxito global
    \end{itemize}
\end{enumerate}

\paragraph{Lista de Verificación para Entrega}

\begin{enumerate}
    \item[$\square$] Rutinas implementadas para las 4 figuras geométricas
    \item[$\square$] Nodo ROS2 funcional con arquitectura publisher/subscriber
    \item[$\square$] Modelo URDF modificado con contenedor, zona de recolección y canecas
    \item[$\square$] Video de funcionamiento en simulación Gazebo (mínimo 2 minutos)
    \item[$\square$] Video de funcionamiento con robot real (mínimo 3 minutos)
    \item[$\square$] Detección visual funcional de 4 figuras con precisión documentada
    \item[$\square$] Diagrama de nodos generado con \texttt{rqt\_graph}
    \item[$\square$] Repositorio Git con código completo y documentación
    \item[$\square$] Plano de planta del montaje físico (vistas superior y lateral)
    \item[$\square$] Diagrama de flujo del sistema en formato Mermaid
    \item[$\square$] Videos con intro institucional del laboratorio
\end{enumerate}

%==============================================================================
\subsection{Criterios de Evaluación}
%==============================================================================

La evaluación del proyecto se realiza mediante rúbrica estructurada en tres componentes con peso equivalente.

\subsubsection{Componente 1: Implementación Técnica (5 puntos)}

\begin{table}[H]
\centering
\caption{Rúbrica - Implementación técnica}
\label{tab:rubrica_tecnica}
\footnotesize
\begin{tabular}{|l|p{10cm}|c|}
\hline
\textbf{Criterio} & \textbf{Descripción} & \textbf{Puntos} \\
\hline
Rutina Cubo & Rutina funcional que transporta cubo desde zona de recolección a caneca roja & 1.0 \\
\hline
Rutina Rectángulo & Rutina funcional que transporta rectángulo a caneca amarilla & 1.0 \\
\hline
Rutina Pentágono & Rutina funcional que transporta pentágono a caneca azul & 1.0 \\
\hline
Rutina Círculo & Rutina funcional que transporta círculo a caneca verde & 1.0 \\
\hline
Simulación Gazebo & Video mostrando operación correcta de las 4 rutinas en entorno simulado & 1.0 \\
\hline
\textbf{Subtotal} & & \textbf{5.0} \\
\hline
\end{tabular}
\end{table}

\subsubsection{Componente 2: Sistema de Percepción (5 puntos)}

\textbf{Variante A - Grupos con Cámara:}

\begin{table}[H]
\centering
\caption{Rúbrica - Visión por computador}
\label{tab:rubrica_vision}
\footnotesize
\begin{tabular}{|l|p{10cm}|c|}
\hline
\textbf{Criterio} & \textbf{Descripción} & \textbf{Puntos} \\
\hline
Nodo ROS2 & Nodo que suscribe a tópico de cámara y publica tipo de figura detectada & 1.5 \\
\hline
Detección de Figuras & Implementación funcional de detección de las 4 figuras mediante OpenCV & 2.0 \\
\hline
Integración & Sistema completo que detecta figura y ejecuta rutina de clasificación correspondiente & 1.5 \\
\hline
\textbf{Subtotal} & & \textbf{5.0} \\
\hline
\end{tabular}
\end{table}

\textbf{Variante B - Grupos con Sistema de Succión:}

\begin{table}[H]
\centering
\caption{Rúbrica - Control de vacío}
\label{tab:rubrica_vacio}
\footnotesize
\begin{tabular}{|l|p{10cm}|c|}
\hline
\textbf{Criterio} & \textbf{Descripción} & \textbf{Puntos} \\
\hline
Nodo Control Teclado & Nodo que escucha teclas y publica comandos para seleccionar figura & 1.5 \\
\hline
Circuito Bomba & Implementación de circuito de control de bomba de vacío con relé & 2.0 \\
\hline
Integración & Sistema que permite seleccionar figura por teclado y ejecutar clasificación & 1.5 \\
\hline
\textbf{Subtotal} & & \textbf{5.0} \\
\hline
\end{tabular}
\end{table}

\subsubsection{Componente 3: Documentación y Evidencias (5 puntos)}

\begin{table}[H]
\centering
\caption{Rúbrica - Documentación}
\label{tab:rubrica_doc}
\footnotesize
\begin{tabular}{|l|p{10cm}|c|}
\hline
\textbf{Criterio} & \textbf{Descripción} & \textbf{Puntos} \\
\hline
Modificación URDF & Contenedor, zona de recolección y canecas agregadas correctamente en Xacro & 1.0 \\
\hline
Video Robot Real & Video en YouTube mostrando clasificación de las 4 figuras con robot físico & 1.5 \\
\hline
Diagramas Técnicos & Diagrama \texttt{rqt\_graph}, plano de planta, diagrama de flujo Mermaid & 1.5 \\
\hline
Repositorio & Código completo, README con instrucciones de uso, estructura organizada & 1.0 \\
\hline
\textbf{Subtotal} & & \textbf{5.0} \\
\hline
\multicolumn{2}{|r|}{\textbf{TOTAL}} & \textbf{15.0} \\
\hline
\end{tabular}
\end{table}

%==============================================================================
\subsection{Recursos de Apoyo}
%==============================================================================

\subsubsection{Documentación Oficial de Referencia}

\begin{itemize}
    \item \textbf{ROS2 Jazzy}: Documentación completa en \url{https://docs.ros.org/en/jazzy/}
    \item \textbf{MoveIt2}: Tutoriales y API en \url{https://moveit.ros.org/}
    \item \textbf{OpenCV}: Documentación de funciones en \url{https://docs.opencv.org/}
    \item \textbf{Gazebo}: Simulación física en \url{https://gazebosim.org/docs}
\end{itemize}

\subsubsection{Comandos de Referencia Rápida}

Los estudiantes deben familiarizarse con los siguientes comandos esenciales de ROS2 para operación y depuración del sistema:

\begin{itemize}
    \item \textbf{Compilación}: Utilizar \texttt{colcon build} con opción de enlaces simbólicos para recompilar workspace tras modificaciones
    \item \textbf{Listar nodos}: Comando \texttt{ros2 node list} muestra todos los nodos activos en el sistema
    \item \textbf{Listar tópicos}: Comando \texttt{ros2 topic list} enumera todos los tópicos disponibles para comunicación
    \item \textbf{Monitorear datos}: Comando \texttt{ros2 topic echo} seguido del nombre de tópico permite observar mensajes en tiempo real (útil para \texttt{/joint\_states})
    \item \textbf{Visualizar arquitectura}: Herramienta \texttt{rqt\_graph} genera diagrama de nodos y tópicos del sistema
    \item \textbf{Árbol de transformadas}: Herramienta \texttt{tf2\_tools view\_frames} genera PDF con jerarquía completa de marcos de referencia
\end{itemize}

Estos comandos constituyen el conjunto mínimo para navegación, depuración y validación del sistema durante desarrollo.

\subsubsection{Herramientas de Depuración}

\begin{table}[H]
\centering
\caption{Herramientas para diagnóstico del sistema}
\label{tab:herramientasDiag}
\small
\begin{tabular}{|l|p{8cm}|}
\hline
\textbf{Herramienta} & \textbf{Uso} \\
\hline
\texttt{rqt\_graph} & Visualización de arquitectura de nodos y tópicos \\
\texttt{rqt\_plot} & Gráficas en tiempo real de datos de tópicos \\
\texttt{rqt\_image\_view} & Visualización de streams de cámara \\
\texttt{tf2\_echo} & Verificación de transformadas entre marcos de referencia \\
\texttt{ros2 bag} & Grabación y reproducción de datos de tópicos \\
\hline
\end{tabular}
\end{table}

%==============================================================================
\subsection{Recomendaciones para Docentes}
%==============================================================================

Basándose en experiencia de implementaciones previas del kit, se recomiendan las siguientes prácticas pedagógicas:

\subsubsection{Organización de Sesiones de Laboratorio}

\begin{enumerate}
    \item \textbf{Sesión 1 - Configuración (4 horas)}
    \begin{itemize}
        \item Instalación guiada de ROS2 y dependencias
        \item Clonación y compilación del workspace
        \item Verificación de funcionamiento básico
        \item Demostración del sistema completo funcionando
    \end{itemize}
    
    \item \textbf{Sesión 2 - Modelado URDF (3 horas)}
    \begin{itemize}
        \item Explicación de formato URDF/Xacro
        \item Modificación guiada del modelo
        \item Ejercicios de agregado de componentes
        \item Validación en RViz
    \end{itemize}
    
    \item \textbf{Sesión 3 - MoveIt2 (4 horas)}
    \begin{itemize}
        \item Conceptos de planificación de movimiento
        \item Uso de Motion Planning Plugin en RViz
        \item Desarrollo de primera rutina de clasificación
        \item Análisis de trayectorias generadas
    \end{itemize}
    
    \item \textbf{Sesión 4 - Programación ROS2 (4 horas)}
    \begin{itemize}
        \item Estructura de nodos en Python
        \item Implementación de publisher/subscriber
        \item Definición de mensajes personalizados
        \item Integración con MoveIt2
    \end{itemize}
    
    \item \textbf{Sesión 5 - Visión por Computador (4 horas)}
    \begin{itemize}
        \item Pipeline de procesamiento de imágenes
        \item Calibración de parámetros HSV
        \item Implementación de detección de figuras
        \item Pruebas con objetos reales
    \end{itemize}
    
    \item \textbf{Sesión 6 - Integración (3 horas)}
    \begin{itemize}
        \item Conexión de todos los componentes
        \item Pruebas de sistema completo
        \item Optimización de parámetros
        \item Documentación final
    \end{itemize}
\end{enumerate}

\subsubsection{Gestión de Grupos de Trabajo}

\begin{itemize}
    \item \textbf{Tamaño de grupos}: 2-3 estudiantes por kit (permite colaboración sin fragmentación excesiva)
    \item \textbf{Asignación de roles}: Rotar responsabilidades entre sesiones (uno programa, otro documenta, otro opera hardware)
    \item \textbf{Control de versiones}: Uso obligatorio de Git con commits documentados
    \item \textbf{Reuniones de seguimiento}: Check-points semanales de 15 minutos por grupo para identificar bloqueos
\end{itemize}

\subsubsection{Puntos Críticos Identificados}

Problemas recurrentes observados en implementaciones previas:

\begin{enumerate}
    \item \textbf{Configuración de entorno}: Errores en paths de ROS2, versiones incompatibles de paquetes
    \begin{itemize}
        \item Solución: Proveer script de instalación automatizado
    \end{itemize}
    
    \item \textbf{Comprensión de transformadas TF}: Dificultad para entender marcos de referencia y conversiones
    \begin{itemize}
        \item Solución: Ejercicio previo de visualización con \texttt{tf2\_tools}
    \end{itemize}
    
    \item \textbf{Calibración de cámara}: Parámetros HSV muy sensibles a iluminación ambiental
    \begin{itemize}
        \item Solución: Proveer rangos iniciales y herramienta interactiva de calibración
    \end{itemize}
    
    \item \textbf{Integración final}: Errores de sincronización entre componentes
    \begin{itemize}
        \item Solución: Enfatizar pruebas modulares antes de integración completa
    \end{itemize}
\end{enumerate}

%==============================================================================
\subsection{Extensiones Avanzadas}
%==============================================================================

Para estudiantes o grupos con progreso acelerado, se sugieren las siguientes extensiones:

\subsubsection{Extensión 1: Clasificación por Múltiples Atributos}

Implementar clasificación basada en combinación de forma y color:
\begin{itemize}
    \item Detección simultánea de forma (círculo, cuadrado, pentágono) y color (rojo, azul, verde, amarillo)
    \item 12 figuras distintas (3 formas × 4 colores)
    \item Matriz de canecas 3×4 para depósito
\end{itemize}

\subsubsection{Extensión 2: Planificación de Secuencia Óptima}

Desarrollo de algoritmo que optimiza orden de clasificación:
\begin{itemize}
    \item Detección de múltiples objetos en zona de recolección
    \item Planificación de secuencia que minimiza distancia total recorrida
    \item Implementación de algoritmo del viajante (TSP) simplificado
\end{itemize}

\subsubsection{Extensión 3: Aprendizaje por Refuerzo}

Entrenamiento de política de agarre mediante RL:
\begin{itemize}
    \item Simulación en Gazebo con variaciones de posición de objeto
    \item Red neuronal que predice pose óptima de agarre
    \item Transferencia de política aprendida a robot real
\end{itemize}

\subsubsection{Extensión 4: Interfaz de Usuario Web}

Desarrollo de dashboard para monitoreo y control:
\begin{itemize}
    \item Visualización en tiempo real de estados del robot
    \item Stream de video de cámara con overlays de detección
    \item Controles para inicio/parada de secuencias
    \item Gráficas de estadísticas de clasificación
\end{itemize}
